\documentclass[12pt,a4paper]{article}

\usepackage[utf8]{inputenc}
\usepackage[T1]{fontenc}
\usepackage[brazil]{babel}
\usepackage{geometry}
\usepackage{xcolor}
\usepackage{listings}
\usepackage{hyperref}
\usepackage[calc]{datetime2}

\geometry{margin=2.5cm}

\definecolor{commentsColor}{rgb}{0.497495, 0.497587, 0.497464}
\definecolor{keywordsColor}{rgb}{0.000000, 0.000000, 0.635294}
\definecolor{stringColor}{rgb}{0.558215, 0.000000, 0.135316}

% Os números ficam na margem interna do quadro, sem ultrapassar sua largura.
\lstdefinestyle{codigo}{
	numbers=left,
	numberstyle=\tiny\color{gray},
	numbersep=6pt,
	xleftmargin=2.5em,
	frame=single,
	framexleftmargin=2em,
	framesep=5pt,
	breaklines=true,
	breakatwhitespace=false,
	columns=fullflexible,
	keepspaces=true,
	basicstyle=\ttfamily\small,
	captionpos=b,
	commentstyle=\color{commentsColor}\textit,
  	keywordstyle=\color{keywordsColor}\bfseries,
  	stringstyle=\color{stringColor},
	postbreak=\mbox{\textcolor{gray}{$\hookrightarrow$}\space},
}
\lstset{style=codigo}

\newcommand{\mescurto}[1]{%
  \ifcase#1\or jan\or fev\or mar\or abr\or mai\or jun\or jul\or ago\or set\or out\or nov\or dez\fi
}

\DTMnewdatestyle{curto}{%
  \renewcommand*{\DTMdisplaydate}[4]{##3 de \mescurto{##2}. de ##1}%
  \renewcommand*{\DTMDisplaydate}{\DTMdisplaydate}%
}

\DTMsetdatestyle{curto}

\DTMsavedate{reportdate}{2026-08-28}

\title{Relatório}
\author{%
	\begin{tabular}{rl}
		Nome: & T.H. de Camargo J. \\
		RA: & 206235
	\end{tabular}%
}

\date{\DTMusedate{reportdate}}

\begin{document}
\maketitle

\section{Introdução}
Apresente aqui o objetivo e o contexto do relatório.

\section{Metodologia}
Descreva os métodos utilizados. Para incluir código com numeração de linhas:

\begin{lstlisting}[language=Python, caption={Exemplo de código}, label={lst:exemplo}]
def soma(a, b):
		return a + b
\end{lstlisting}

Consulte o Código~\ref{lst:exemplo} ao discutir a implementação.
Também é possível incluir um arquivo externo com

\lstinputlisting[language=c, caption={Something}, label={label}]{../a.c}

\section{Resultados e discussão}
Apresente e discuta os resultados.

\section{Conclusão}
Resuma as conclusões do trabalho.

\end{document}
